\documentclass[tikz,border=6pt]{standalone}
\usepackage{fontspec}
\usepackage{xeCJK}
\setmainfont{Times New Roman}
\setCJKmainfont{Hiragino Sans GB}
\setCJKsansfont{Microsoft YaHei}
\setCJKmonofont{STSong}
\usepackage{xcolor}

\definecolor{blockblue}{HTML}{4A90D9}

\begin{document}
\begin{tikzpicture}[line width=1.2pt]
  \node[
    draw=blockblue,
    rounded corners=3pt,
    inner sep=8pt,
    align=left
  ] at (0,0) {
    \begin{tabular}{p{3.8cm}p{9.0cm}}
      \textbf{由频率响应整定得到} & $K_c=0.7,\quad T_c=0.2566$ \\
      比例控制器 & $K_P=0.5K_c=0.35$ \\
      比例积分控制器 & $K_P=0.4K_c=0.28,\quad T_I=0.8T_c\approx0.21$ \\
      比例积分微分控制器 & $K_P=0.6K_c=0.42,\quad T_I=0.5T_c\approx0.13,\quad T_D=0.125T_c\approx0.03$ \\
      说明 & 当闭环不稳定时，虽然阶跃响应法不可用，但频率响应整定仍然可用 \\
    \end{tabular}
  };
\end{tikzpicture}
\end{document}
